\section{Simulation-methods comparison with OpenRocket}\label{app:or-methods}

Where \cref{app:or-features} compares the two simulators feature by feature, this appendix
compares them \emph{method by method}: for each area of flight physics, what OpenRocket
computes, what QtRocket computes, and a short assessment of the distance between the two.
It is written for the reader who wants the actual equations. OpenRocket's methods are taken
from its technical documentation\footnote{Sampo Niskanen, \emph{OpenRocket technical
documentation, for OpenRocket version 13.05}, 2013-05-10,
\url{https://openrocket.sourceforge.net/techdoc.pdf} --- the updated form of Niskanen's 2009
Master's thesis. Cited hereafter as ``techdoc'' with its own section and equation numbers.}
and, where the 2013 techdoc has been overtaken by later development (geodesy, Coriolis,
multi-level wind), from the current source.\footnote{OpenRocket source, \code{unstable}
branch of \url{https://github.com/openrocket/openrocket} (fetched 2026-07-07), files under
\code{core/src/main/java/info/openrocket/core/}. Cited hereafter by class name.}
As \cref{app:or-context} establishes, OpenRocket-side claims are condensed from those
sources --- researched, but not line-verified in the way QtRocket claims are. Every QtRocket
claim cites the pinned commit \code{254cd41}. \Cref{tab:or-methods-summary} condenses the
whole appendix into one table.

\subsection{Flight dynamics}\label{app:or-methods-dynamics}

\emph{OpenRocket.} Ascent is a genuine six-degree-of-freedom (6-DOF) simulation: world-frame
position and velocity plus a unit orientation quaternion and body angular velocity, written
as a first-order system and stepped together (techdoc \S 4.2.1, \S 4.2.4, eqs.~4.17--4.19).
Every force evaluation recomputes local wind, angle of attack (AoA), and Mach number, then
the full aerodynamic coefficient set at those instantaneous conditions, then thrust,
gravity, mass, and moments of inertia (techdoc \S 4.2.4, steps 0--6). In the code the linear
equation is assembled in rocket coordinates as $(-f_N,\,-f_{side},\,T-D_{axial})/m$, rotated
to the world frame, with gravity and Coriolis acceleration added; the pitch, yaw, and roll
moments are referred to the center of gravity (CG) and divided by exactly \emph{two} scalar
moments of inertia --- longitudinal for pitch/yaw, axial for roll --- so the inertia tensor
is treated as diagonal and axisymmetric, never as a full $3\times 3$
tensor.\footnote{\code{simulation/RK4SimulationStepper.java},
\code{computeAcceleration()}.} A deliberate $\pm 0.0005$ uniform random perturbation is
added to $C_m$ and $C_{yaw}$ at every evaluation ``to prevent over-perfect
flight''.\footnote{\code{RK4SimulationStepper}, constant \code{PITCH\_YAW\_RANDOM}.}
The 6-DOF model is also deliberately abandoned when it stops earning its keep: ``When a
parachute is deployed the rocket typically splits in half, and it is no longer practical to
compute the orientation of the rocket'' (techdoc \S 4.2.5), so recovery descent and the
descent of unstable, tumbling stages (techdoc \S 3.5) run as 3-DOF point masses under a
different stepper. Fidelity is matched to flight phase.

\emph{QtRocket.} The integrated state is three linear degrees of freedom, full stop. The
entire physics reaching the integrator is one lambda,
$\dot{\mathtt{state}} = \mathtt{rate}$,
$\dot{\mathtt{rate}} = \mathtt{getForces}(t)/\mathtt{getMass}(t)$
\src{sim/Propagator.cpp}{36}, with thrust hardwired along world $+z$
\src{model/RocketModel.cpp}{71} --- there is no attitude for it to follow.
\code{StateData} reserves quaternion orientation and rate, a direction-cosine matrix, and
3-2-1 Euler-angle slots, and nothing anywhere writes them \src{sim/StateData.h}{46}; the
orientation integrator exists as a commented-out member with a design note
\src{sim/Propagator.h}{103}; \code{getTorques()} is declared on the bridge interface
\src{model/Propagatable.h}{30}, implemented as the zero vector
\src{model/RocketModel.cpp}{94}, and called by no production code. All of this is dormant
seam, inventoried in \cref{sec:incomplete} and described in \cref{sec:simcore}.

\begin{keyinsight}[The inertia inversion]
The data models run in the opposite direction from the dynamics. OpenRocket integrates real
attitude but represents inertia as two scalars; QtRocket integrates no attitude but records
the \emph{full} $3\times 3$ composite inertia tensor about the instantaneous composite CM
into every retained state sample \src{sim/StateData.h}{62}, refreshed each step through
\code{writeMassProperties} \src{sim/Propagator.cpp}{132} and pinned by a test that watches
$I_{xx}$ shrink through the burn and freeze after burnout
\src{tests/PhysicsIntegrationTests.cpp}{356}. QtRocket thus carries a strictly richer
rotational data model than the one OpenRocket actually flies with --- and integrates none of
it. The recorded tensor is the declared seam the future 6-DOF rotational ODE will read.
\end{keyinsight}

The distance here is one of kind, not tuning: for weathercocking, gravity turns, spin, or
dispersion, OpenRocket computes and QtRocket cannot yet.

\subsection{Aerodynamic coefficients}\label{app:or-methods-aero}

\emph{OpenRocket.} The normal-force machinery is an \emph{extended} Barrowman method,
re-evaluated at the instantaneous flight condition on every force call. Instead of the
classic $\alpha \to 0$ derivative, OpenRocket defines, for $\alpha > 0$,
\begin{equation}
C_{N_\alpha} \equiv \frac{C_N}{\alpha}, \qquad C_{m_\alpha} \equiv \frac{C_m}{\alpha}
\end{equation}
so the whole coefficient set depends on the actual AoA and Mach number (techdoc \S 3.1.1,
eq.~3.8). The extensions beyond classic Barrowman, each with its techdoc equation: body lift
per Galejs, $C_N^{\text{lift}} = K\,(A_{\mathrm{plan}}/A_{\mathrm{ref}})\sin^2\alpha$ with
$K \approx 1.1$, acting at the planform centroid (eqs.~3.26--3.27); free-form fin planforms
via numerically integrated mean-aerodynamic-chord (MAC) integrals (eqs.~3.30--3.32);
arbitrary fin counts with empirical interference multipliers falling to $0.750$ beyond eight
fins (eq.~3.54); fin--body interference $K_{T(B)} = 1 + r_t/(s + r_t)$ (eqs.~3.55--3.56);
pitch and yaw damping proportional to $\omega^2/v_0^2$ (eqs.~3.59--3.60); and full roll
forcing/damping strip theory for canted fins (eqs.~3.66, 3.69--3.72). Compressibility runs
through everything: the Prandtl factor $1/\beta$ with $\beta = \sqrt{|M^2-1|}$
(eqs.~3.13--3.14), Busemann third-order strips for supersonic fins (eqs.~3.41--3.49), and a
moving fin center of pressure (CP): fixed at quarter-MAC below $M\,0.5$, the empirical
$X_f/\bar c = (\mathcal{AR}\beta - 0.67)/(2\mathcal{AR}\beta - 1)$ above $M\,2$, and a
fifth-order polynomial in $M$ with matched values and derivatives bridging the transonic gap
between them (eqs.~3.35--3.36).

\emph{QtRocket.} The Barrowman layer is the classic zero-AoA linearization, computed per
part and folded over the placement-resolved tree (\cref{sec:aero}). A nose cone contributes
$\mathtt{cnAlpha} = 2\pi\,\mathtt{baseRadius}^2/\mathtt{refArea}$ with its CP at $2L/3$ aft
of the tip \src{model/parts/ConicalNoseCone.cpp}{64}; a fin set contributes the standard
trapezoidal-fin slope \src{model/parts/FinSet.cpp}{77} multiplied by an interference factor
\emph{identical in form to OpenRocket's} eq.~3.55,
$K = 1 + \mathtt{bodyRadius}/(\mathtt{span} + \mathtt{bodyRadius})$
\src{model/parts/FinSet.cpp}{76}; body tubes contribute zero slope, exactly as slender-body
theory demands. Contributions are summed as slope-weighted CP moments so the composite CP is
a guarded ratio \src{model/parts/Part.cpp}{231}, \src{sim/Aero.h}{51}. What the layer does
\emph{not} have: any $\alpha$ dependence beyond the linear slope, any body-lift term, any
damping or roll coefficient, and any Mach dependence at all --- nothing in the tree computes
a Mach number, and the atmosphere's speed-of-sound accessor has no production caller
\src{sim/USStandardAtmosphere.cpp}{143}. And the layer is dormant: the test suite states
``nothing in production consumes this'' verbatim \src{model/tests/AeroTests.cpp}{18}.

The assessment writes itself: formula by formula the two projects share an ancestor, and
where they overlap (fin slope, interference factor, CP-moment composition) they agree.
OpenRocket then evaluates its extended version at flight conditions every step and flies on
the result; QtRocket's classic version is evaluated at $\alpha \to 0$, with no Mach anywhere,
by tests alone.

\subsection{Drag}\label{app:or-methods-drag}

\emph{OpenRocket.} Zero-AoA drag is a per-component, semi-empirical build-up (techdoc
\S 3.4): skin friction, body and fin pressure drag, base drag, and parasitic (launch-guide)
drag, with interference drag explicitly neglected. Skin friction assumes a fully turbulent
boundary layer, using the smooth-wall correlation and a roughness-limited floor
\begin{equation}
C_f = \frac{1}{(1.50\,\ln R - 5.6)^2}
\qquad\text{and}\qquad
C_f = 0.032\left(\frac{R_s}{L}\right)^{0.2}
\end{equation}
selected piecewise by a roughness-critical Reynolds number, with separate subsonic and
supersonic compressibility corrections (eqs.~3.78--3.84) and body/fin geometry factors when
scaled to the reference area (eq.~3.85). Nose pressure drag is per-shape: a stagnation-based
blunt-cylinder limit, a closed form for cones, an ogive shape-factor scaling, and stored
experimental transonic curves for ellipsoid/power/parabolic/Haack profiles, interpolated in
Mach and fineness ratio (techdoc App.~B); high-subsonic values are bridged by a constrained
power law (eq.~3.87). Base drag is empirical,
\begin{equation}
(C_{D\bullet})_{\mathrm{base}} =
\begin{cases}
0.12 + 0.13\,M^2 & M < 1\\[2pt]
0.25/M & M > 1
\end{cases}
\end{equation}
with a thrusting motor's cross-section subtracted from the base area (eq.~3.94); launch lugs
are blunt cylinders at stagnation pressure with a flow-through correction (eqs.~3.95--3.96).
Finally the zero-AoA total (eq.~3.97) is stretched over AoA by an empirical two-part
polynomial $C_A(\alpha) = C_{D_0} f(\alpha)$ with $f(0)=1$, a hump of $1.3$ at
$\alpha = 17^\circ$, and $f(90^\circ) = 0$ (techdoc \S 3.4.7).

\emph{QtRocket.} The production drag model is one line
\src{model/RocketModel.cpp}{88}:
\begin{equation}
\vec F_{\mathrm{drag}} = -\tfrac{1}{2}\,\rho(\mathtt{altitude})\,
\lVert \mathtt{velocity} \rVert \,
\mathtt{dragCoefficient}\cdot\mathtt{referenceArea}\cdot\mathtt{velocity},
\end{equation}
a single constant $C_d$ (default $1.0$, \src{model/RocketModel.h}{138}) times a single
constant reference area (default $1.134\times 10^{-3}\,\mathrm{m}^2$, a 38\,mm disc,
\src{model/RocketModel.h}{141}), opposing the full velocity vector, with density sampled
at the trial altitude clamped to $z \ge 0$ \src{model/RocketModel.cpp}{85}. There is no
Mach or Reynolds dependence, no AoA dependence, and no per-component build-up: the
\code{AeroComponent} value type carries a per-part \code{cd} field for exactly that purpose,
and every concrete part returns zero --- the body tube's comment defers skin friction over
its (unconsumed) wetted area to future work \src{model/parts/BodyTube.cpp}{41},
\src{model/parts/ConicalNoseCone.cpp}{67}, \src{model/parts/FinSet.cpp}{89},
\src{model/parts/BodyTube.h}{51}.

\begin{hardfail}[The reference area never reaches the drag law]
QtRocket implements the max-frontal-disc reference-area convention correctly ---
\code{deriveReferenceAreaFromGeometry()} \src{model/RocketModel.cpp}{48} walks the part tree
taking the maximum frontal disc \src{model/parts/Part.cpp}{251} --- but the function has no
production caller; its only call site outside its own definition is a test
\src{model/tests/AeroTests.cpp}{201}. Design load applies a saved area only when the design
was saved with a manual override \src{model/DesignSerializer.cpp}{263}, and all 25 bundled
design fixtures carry \code{referenceAreaOverridden="false"}, so every loaded design flies
on whatever scalar was previously staged --- by default the 38\,mm disc --- regardless of the
airframe's real diameter. Drag does not scale with airframe class at this commit. \fail
\end{hardfail}

Distance: OpenRocket's drag is five semi-empirical families evaluated per component per
step; QtRocket's is one user-supplied number, currently wired to one defective constant.
This is the widest methods gap in the whole comparison.

\subsection{Atmosphere and wind}\label{app:or-methods-atmosphere}

Both atmospheres are members of the same family --- the International Standard Atmosphere
(ISA) and the US Standard Atmosphere 1976 are identical below 32\,km --- and this is the one
area where the two simulators are nearly peers. OpenRocket linearly interpolates temperature
between eight ISA layer bases (0 to 84.852\,km) and integrates the hydrostatic/ideal-gas
pair within each layer (techdoc \S 4.1.1, Table~4.1); the current code adds a launch-site
override that inserts an extra layer with a back-computed sea-level lapse rate, converts
geometric to geopotential height with $R = 6\,356\,766$\,m, and precomputes layer conditions
for speed.\footnote{\code{models/atmosphere/ExtendedISAModel.java}.} QtRocket's
\code{USStandardAtmosphere} hard-codes the same first seven layer anchors (0 to 71\,km ---
it has no 84.852\,km top row) with published base pressures and densities
\src{sim/USStandardAtmosphere.cpp}{78}, and evaluates the closed-form gradient-layer power
law and isothermal exponential exactly within each layer
\src{sim/USStandardAtmosphere.cpp}{121}, the density exponent carrying the ideal-gas
$-1$ \src{sim/USStandardAtmosphere.cpp}{107}. The differences are at the edges: QtRocket has
no launch-site override (the registry offers exactly Constant, US Standard 1976, and Vacuum,
\src{sim/Environment.h}{74}), no geopotential conversion, throws below $0$\,m
\src{utils/Bin.cpp}{52} (its one consumer clamps, \src{model/RocketModel.cpp}{85}), and
extrapolates the 71\,km layer upward without bound \src{utils/Bin.cpp}{59}. Against that,
QtRocket's implementation is the more aggressively \emph{verified} of the two: pressure and
temperature match the bundled NOAA report tables to 0.1\% relative at twenty altitudes from
0 to 80\,km \src{sim/tests/USStandardAtmosphereTests.cpp}{36},
\src{sim/tests/USStandardAtmosphereTests.cpp}{64}, and density to 0.1\%/0.5\%
\src{sim/tests/USStandardAtmosphereTests.cpp}{10},
\src{sim/tests/USStandardAtmosphereTests.cpp}{24}.

Wind is not close. OpenRocket models wind as a steady average plus pink-noise turbulence
with spectrum $1/f^{\alpha}$, $\alpha = 5/3$, generated by a two-pole Kasdin filter over
Gaussian white noise at a fixed 20\,Hz and scaled by the filter's precomputed standard
deviation 2.252 (techdoc \S 4.1.2, eqs.~4.5--4.6);\footnote{Constants verbatim in
\code{models/wind/PinkNoiseWindModel.java}: \code{ALPHA} $=5/3$, \code{POLES} $=2$,
\code{STDDEV} $=2.252$, \code{DELTA\_T} $=0.05$.} since release 24.12, arbitrary altitude
levels each carry their own pink-noise model, with the instantaneous wind vectors of the two
adjacent levels linearly interpolated component-wise.\footnote{%
\code{models/wind/MultiLevelPinkNoiseWindModel.java}; OpenRocket ReleaseNotes.md, 24.12.}
QtRocket's \code{WindModel} is a concrete class returning the zero vector
\src{sim/WindModel.cpp}{16} with zero call sites and no \code{Environment} slot
\src{sim/Environment.h}{103}; drag is computed from inertial velocity, so the airmass is
implicitly at rest \src{model/RocketModel.cpp}{88}.

\subsection{Gravity and geodesy}\label{app:or-methods-gravity}

The 2013 techdoc states flat-Earth Cartesian coordinates and ``the coriolis effect \ldots{}
is not simulated'' (techdoc \S 4.2) --- and is obsolete on this point. Current OpenRocket
offers three geodetic strategies --- flat, spherical great-circle (the default), and a full
Vincenty direct geodesic on the WGS84 (World Geodetic System 1984) ellipsoid --- and adds
Coriolis acceleration $-2\boldsymbol{\Omega}\times\vec v$ to every force evaluation under
the latter two.\footnote{\code{util/GeodeticComputationStrategy.java};
\code{simulation/RK4SimulationStepper.java} and \code{simulation/AbstractEulerStepper.java}.}
Gravity is either a user constant or Somigliana normal gravity on the ellipsoid,
\begin{equation}
g_0(\lambda) = 9.7803267714\;
\frac{1 + 0.00193185138639\,\sin^2\lambda}{\sqrt{1 - 0.00669437999013\,\sin^2\lambda}},
\end{equation}
with an inverse-square altitude correction
$g = g_0\,\bigl(R_E/(R_E+h)\bigr)^2$.\footnote{\code{models/gravity/WGSGravityModel.java}.}

QtRocket offers two gravity models (\cref{sec:environment}): a constant field
$(0,0,-g_0)$ with $g_0 = 9.80665\,\mathrm{m/s^2}$ \src{sim/ConstantGravityModel.h}{18},
\src{utils/math/Constants.h}{10} (the default, \src{sim/Environment.h}{33}), and a true
Newtonian central force over a spherical Earth:
\begin{equation}
\vec a = -\frac{\mathtt{earthGM}}{r^3}\,\vec r_{\mathrm{geo}},
\qquad \vec r_{\mathrm{geo}} = (x,\ y,\ z + \mathtt{groundLevel}),
\end{equation}
where the pad's geocentric radius comes from a geoid model
\src{sim/SphericalGravityModel.cpp}{33}, \src{sim/SphericalGravityModel.cpp}{36}. The only
geoid returns the WGS84 mean radius $6\,371\,008.8$\,m at every latitude and longitude
\src{sim/SphericalGeoidModel.cpp}{18}, and the gravity model caches it at the placeholder
site $(0,0)$ \src{sim/SphericalGravityModel.cpp}{20}. The spherical model is in one sense
the more literal of the two central-force treatments --- a downrange displacement picks up a
genuine inward horizontal component, pinned by test
\src{sim/tests/SphericalGeoidAndGravityTests.cpp}{82}, and the surface magnitude matches
$GM/R^2$ to $10^{-12}$ relative \src{sim/tests/EnvironmentTests.cpp}{123} --- but QtRocket
has no latitude dependence, no Coriolis term, and no launch-site concept at all. For
low-altitude model rockets both simulators' options agree to well under 1\%; the OpenRocket
extras matter for long boosted flights and dispersion analysis, which QtRocket does not yet
attempt.

\subsection{Numerical integration}\label{app:or-methods-integration}

\begin{physbox}[What an embedded pair buys]
A Runge--Kutta step of order $p$ commits an unknown local error of order $h^{p+1}$. An
\emph{embedded pair} evaluates one shared set of stage derivatives and combines them twice,
with weight sets of order $p$ and $p+1$; the difference of the two results is a nearly free
estimate of the lower-order solution's local error, which a controller can compare against a
tolerance to \emph{reject and retry} a too-large step or grow a wastefully small one. A
fixed-step or heuristically stepped integrator, however well tuned, never measures its own
error; an embedded-pair integrator measures it every step. The cost is extra stages (six for
Fehlberg's 4(5) pair versus four for classical RK4) and the machinery of acceptance and
rejection.
\end{physbox}

\emph{OpenRocket.} The ascent stepper is classical fixed-formula RK4 (techdoc \S 4.2.4,
eqs.~4.20--4.21) with the quaternion updated multiplicatively and renormalized each step.
Step-size control is heuristic and a priori: after evaluating only $k_1$, the stepper takes
the minimum of eight candidate steps --- the user step (default 0.05\,s, divided by 5 on the
launch rod), the time to the next scheduled flight event, a $3^\circ$ default cap on pitch
angle per step, a roll-angle step of $2\times 28.32^\circ$ (deliberately an uneven division
of the circle, to sample the wind), a $2^\circ$ roll-rate-change cap, a $4^\circ$
pitch/yaw-acceleration cap, a launch-rod resolution term, and a $1.5\times$ growth limiter
--- floored at 0.001\,s.\footnote{\code{simulation/RK4SimulationStepper.java},
\code{step()}; the candidate list is verbatim from the code comment. Defaults:
\code{RECOMMENDED\_TIME\_STEP} $=0.05$\,s, \code{RECOMMENDED\_ANGLE\_STEP} $=3^\circ$,
\code{MIN\_TIME\_STEP} $=0.001$\,s.} There is no error estimate and no step rejection: the
$k_1$ rates decide the step before it is taken, and the step is never re-integrated. (An RK6
stepper with the same heuristic selection exists on the development branch; it is not in the
24.12 release notes and should be treated as unreleased.) Recovery and tumble descend under
a separate 3-DOF second-order Euler stepper with a 0.5\,s base step, exact root-solves to
land precisely on apogee and ground crossings, and a jerk-based oscillation
guard.\footnote{\code{simulation/AbstractEulerStepper.java}.}

\emph{QtRocket.} The \code{Integrator} facade offers two runtime-selectable backends
\src{sim/Integrator.h}{32}: a fixed-step classical RK4 \src{sim/RK4Solver.h}{27} and a true
Runge--Kutta--Fehlberg 4(5) embedded pair whose Butcher tableau matches the textbook
Fehlberg coefficients digit for digit \src{sim/RK45Solver.h}{172}. Each attempt evaluates
six stages at their proper node times, forms both the fourth- and fifth-order estimates, and
accepts or rejects on the estimated local error (\cref{lst:or-methods-rkf45}): rejection
shrinks $h$ by $(\mathtt{tol}/\mathtt{err})^{1/4}$ floored at $0.1$ and \emph{retries the
step}; acceptance grows the next guess by $(\mathtt{tol}/\mathtt{err})^{1/5}$ capped at
$5\times$ \src{sim/RK45Solver.h}{141}, clamps it to $h_{\max} = 10\times$ the seeded step
\src{sim/RK45Solver.h}{66}, \src{sim/RK45Solver.h}{147} --- the clamp is a correctness
requirement, since exactly-integrated coasting drives $\mathtt{err}\to 0$ and unbounded
growth --- and advances with the fifth-order estimate. Below
$h_{\min} = 10^{-15}$\,s the solver throws \src{sim/RK45Solver.h}{80}, which the run loop
converts to a typed \code{IntegratorError} verdict \src{sim/Propagator.cpp}{171}
(\cref{sec:simcore}).

\begin{listing}[htbp]
\begin{minted}{cpp}
         // Local error estimate: magnitude of the 5th-vs-4th-order difference across state and rate.
         const double err = std::sqrt((y5State - y4State).squaredNorm()
                                      + (y5Rate - y4Rate).squaredNorm());

         if(err > tol)
         {
            // Reject: shrink the step and retry, flooring the shrink factor at 0.1.
            const double scale = std::pow(tol / err, 0.25);
            h *= (scale < 0.1) ? 0.1 : scale;
            continue;
         }
\end{minted}
\caption{The step-rejection core OpenRocket's ascent stepper has no analogue of: an
embedded-pair error estimate compared against a tolerance, with the step retried when it
fails. \srccap{sim/RK45Solver.h}{120}}
\label{lst:or-methods-rkf45}
\end{listing}

\begin{keyinsight}[What each scheme does and does not control]
OpenRocket's heuristics bound the \emph{geometry} of a step --- degrees of pitch or roll per
step --- which is an excellent proxy for the stiff part of a 6-DOF problem, but they never
measure integration error and cannot detect a step that was numerically too coarse; accuracy
rests on the caps being conservative. QtRocket's RKF45 measures the local truncation error
of exactly what it integrates and rejects steps that exceed tolerance --- a genuine
methodological advantage, and the one area where QtRocket's numerics are strictly ahead ---
but three qualifiers bound the claim. The error norm is a single absolute Euclidean norm
mixing meters and meters/second \src{sim/RK45Solver.h}{121}, flagged in-source as a possible
refinement \src{sim/RK45Solver.h}{36}; there is no safety factor or PI controller, only the
$0.1/5\times$ clamps; and the tolerance (default $10^{-6}$, \src{sim/RK45Solver.h}{46}) is
reachable from no front end --- \code{setErrorTolerance}/\code{setMaxStepSize} have no
production callers \src{sim/tests/RK45SolverTests.cpp}{112}. And with only three linear DOF
integrated, the adaptive machinery currently controls a far easier problem than OpenRocket's.
\end{keyinsight}

Head to head on the same 3-DOF flight, QtRocket's RKF45 reproduces its RK4 apogee to about
$0.05\%$ while taking roughly $530$ steps where RK4 at the same seed takes about $4050$
\src{tests/PhysicsIntegrationTests.cpp}{164}, \src{tests/PhysicsIntegrationTests.cpp}{171}.
QtRocket's default fixed step is 0.01\,s \src{sim/Propagator.h}{110} against OpenRocket's
default 0.05\,s --- the finer default compensating for the absence of angle-based caps.

\subsection{Motors}\label{app:or-methods-motors}

Both simulators represent thrust identically in kind: a measured thrust curve, linearly
interpolated between samples, zero outside the burn. OpenRocket's
\code{ThrustCurveMotor} carries (time, thrust, CG$+$mass) sample lists and linearly
interpolates all three;\footnote{\code{motor/ThrustCurveMotor.java}.} QtRocket's
\code{ThrustCurve} interpolates thrust the same way \src{model/ThrustCurve.cpp}{83},
including a regression-pinned ramp from an implicit $(0,0)$ origin when a curve omits its
first sample \src{model/ThrustCurve.cpp}{76} (\cref{sec:motors}).

Mass depletion is where they part. When a RockSim \code{.rse} file supplies measured
per-sample mass and CG, OpenRocket uses those samples directly; only for RASP \code{.eng}
files (which carry just total and propellant mass) does it fall back to distributing mass
loss proportional to burned impulse,
$\Delta m_i \propto \tfrac12 (f_i + f_{i+1})\,\Delta t_i$, scaled to consume exactly the
propellant mass.\footnote{\code{file/motor/AbstractMotorLoader.calculateMass()}.} QtRocket
\emph{always} reconstructs mass under the constant-Isp assumption, even when measured data
is present: the RSE loader reads only each sample's \code{t} and \code{f} attributes and
discards the file's per-sample \code{m} and \code{cg} \src{model/RSEDatabaseLoader.cpp}{88},
and \code{computeMassCurve} builds a 128-point table \src{model/MotorModel.cpp}{134} by
left-Riemann depletion proportional to delivered impulse,
\begin{equation}
\mathtt{propMass}_{i+1} = \mathtt{propMass}_i
- F(t_i)\,\Delta t\,\frac{\mathtt{propWeight}}{\mathtt{totalImpulse}},
\qquad \Delta t = \frac{\mathtt{burnTime}}{127},
\end{equation}
\src{model/MotorModel.cpp}{141}, forcing the final point to exact depletion so post-burn
mass is exactly the casing mass \src{model/MotorModel.cpp}{143}. The physical content is the
same as OpenRocket's RASP fallback --- mass flow proportional to thrust --- so the practical
difference appears only for motors whose measured burn deviates from constant effective
exhaust velocity. Two QtRocket-side quirks are worth restating from \cref{sec:motors}: mass
and thrust gate burnout on different clocks (metadata \code{burnTime} versus the curve's
last sample time, \src{model/MotorModel.cpp}{36}, \src{model/MotorModel.cpp}{72}), and a
derived \code{isp} member is computed and never read \src{model/MotorModel.cpp}{130}.

Ignition and delays: OpenRocket treats motor ignition, burnout, and ejection charges as
flight events, each trigger optionally carrying a delay time --- this is how staging and
dual deployment are expressed (techdoc \S 4.2.6, Table~4.2). QtRocket ignites its single
motor exactly once, at $t = 0$, from \code{RocketModel::launch()}
\src{model/RocketModel.cpp}{105}; \code{ThrustCurve::setIgnitionTime} has zero call sites
\src{model/ThrustCurve.h}{31}, and ejection-delay lists are parsed by every loader and
consumed by nothing \src{model/MotorModel.h}{303}.

\subsection{Recovery physics}\label{app:or-methods-recovery}

OpenRocket has a complete recovery subsystem. A parachute is a drag device with default
$C_D = 0.8$ on its canopy area;\footnote{\code{rocketcomponent/Parachute.java},
\code{DEFAULT\_CD}; techdoc \S 4.2.5.} a streamer's drag comes from Niskanen's own
wind-tunnel experiments,
\begin{equation}
C_{D} = 0.034\cdot
\frac{\rho_m + 25\ \mathrm{g/m^2}}{105\ \mathrm{g/m^2}}\cdot
\frac{l + 1\ \mathrm{m}}{l},
\end{equation}
good to roughly 15\% in descent rate (techdoc App.~C, eq.~C.6); tumbling stages use a
two-component drop-test model with $C_{D,f} = 1.42$ and $C_{D,bt} = 0.56$ (techdoc \S 3.5,
eq.~3.99). Deployment is driven by the flight-event system --- apogee detection,
altitude-threshold main deployment, ejection-charge delays --- and switches the active
stepper from 6-DOF RK4 to the 3-DOF Euler landing stepper (techdoc \S 4.2.5--4.2.6).

QtRocket has none of this: no recovery device exists in the part vocabulary, no deployment
event exists in the simulation loop, and descent is simply the same constant-$C_d$ point
mass falling until the nominal termination condition --- descending \emph{and} below the
launch site, $z < 0 \wedge v_z < 0$ \src{model/RocketModel.cpp}{65} --- ends the flight. The
only event-like machinery is the typed termination taxonomy itself
\src{sim/Propagator.h}{35}. Recovery is an absence, not a dormancy; \cref{sec:incomplete}
carries it in the consolidated inventory.

\subsection{Validation posture}\label{app:or-methods-validation}

OpenRocket published flight-data validation. Against altimeter flights of a small model
rocket, apogee error was $+16\%$ and $+7\%$ (versus RockSim's $+24\%$ and $+19\%$ on the
same flights); a hybrid-powered flight came in at $-16\%$; a magnetometer-instrumented
rolling rocket spun at 16\,rev/s where the simulation predicted about 8 --- roll rate
under-predicted by roughly $2\times$, cause unknown; wind-tunnel comparison put the
simulated CP within a few percent of body length at subsonic speeds. The headline claim is
CP and drag-coefficient accuracy of ``approximately 10\% at subsonic velocities'' (techdoc
ch.~6--7).

QtRocket, at this commit, has no whole-flight field validation: no test compares a simulated
trajectory against a recorded flight, and every oracle in the six test binaries is analytic
or table-derived. What it has instead is unusually dense component-level
\emph{verification}: atmosphere properties to 0.1--0.5\% of the NOAA report tables
\src{sim/tests/USStandardAtmosphereTests.cpp}{36}; the RKF45 solver against exact solutions
--- constant acceleration to $10^{-9}$ \src{sim/tests/RK45SolverTests.cpp}{43}, a harmonic
oscillator to $10^{-4}$ \src{sim/tests/RK45SolverTests.cpp}{126}, exponential decay to
$10^{-5}$ \src{sim/tests/RK45SolverTests.cpp}{150}, a thrust-ramp analogue to $10^{-6}$
\src{sim/tests/RK45SolverTests.cpp}{175}; a ballistic arc's apogee against $v_0^2/2g$ to
0.5\% \src{tests/PropagatorTests.cpp}{178}; the drag law against the analytic
terminal-velocity balance to $10^{-6}$\,N \src{tests/PhysicsIntegrationTests.cpp}{262};
inverse-square surface gravity to $10^{-12}$ relative
\src{sim/tests/EnvironmentTests.cpp}{123}.

\begin{keyinsight}[Verification is not validation]
The two projects' accuracy numbers answer different questions and should not be compared
digit for digit. OpenRocket's $\sim$10\% is \emph{validation}: model versus nature, with all
modeling error included --- and its published misses (roll off $2\times$, supersonic drag
over-predicted) are part of the claim's honesty. QtRocket's $10^{-4}$-to-$10^{-12}$ oracles
are \emph{verification}: code versus its own stated mathematics, guaranteeing the
implementation computes the declared model correctly while saying nothing about whether that
model matches a real flight. Given the force model this appendix has documented --- constant
$C_d$ on a defective reference area, no wind, no attitude --- QtRocket could not currently
support an OpenRocket-style 10\% whole-flight claim, and it does not make one. The honest
statement is: OpenRocket's physics is validated to $\sim$10\%; QtRocket's physics is
verified to compute exactly what it says, and what it says is still a point mass.
\end{keyinsight}

\subsection{Summary}\label{app:or-methods-summary-sec}

\begin{table}[htbp]
\small
\caption{Methods at a glance. OpenRocket per its techdoc and current source; QtRocket per
the pinned commit \code{254cd41}. ``Dormant'' marks machinery that is built and tested with
zero production consumers.}
\label{tab:or-methods-summary}
\begin{tabularx}{\textwidth}{@{}>{\raggedright\arraybackslash}p{2.6cm}XX@{}}
\toprule
\textbf{Area} & \textbf{OpenRocket} & \textbf{QtRocket} \\
\midrule
Degrees of freedom & 6-DOF quaternion ascent, two scalar inertias; deliberate 3-DOF
degradation for recovery and tumble & 3-DOF point mass; quaternion slots and full
$3\times 3$ recorded inertia dormant \\
\addlinespace
Aero coefficients & Extended Barrowman at instantaneous AoA/Mach each step: body lift,
interference, damping, roll, transonic fin CP & Classic Barrowman at $\alpha\to 0$, no Mach
dependence; complete, tested, dormant \\
\addlinespace
Drag & Per-component semi-empirical build-up: friction with roughness floor, per-shape wave
drag, base $0.12+0.13M^2$, lug, $C_A(\alpha)$ hump & Single constant $C_d \times$ constant
reference area; geometry-derived area implemented but never called \\
\addlinespace
Atmosphere & ISA, 8 layers, launch-site override, geopotential conversion & US Standard
1976, 7 layers, closed-form; no overrides; verified 0.1--0.5\% to 80\,km \\
\addlinespace
Wind & Average $+$ 2-pole pink-noise turbulence, 20\,Hz; multi-level profiles (24.12) &
Zero-wind stub, zero call sites \\
\addlinespace
Gravity/geodesy & Somigliana latitude gravity, flat/spherical/WGS84-Vincenty positions,
Coriolis & Constant $g_0$ or true inverse-square over mean-radius geoid; no latitude, no
Coriolis, no launch site \\
\addlinespace
Integration (ascent) & RK4, heuristic a-priori step caps (8 candidates, $3^\circ$ angle), no
error estimate, no rejection & RK4 or embedded-pair RKF45 with acceptance/rejection,
$h_{\max}$ clamp; tolerance not user-reachable \\
\addlinespace
Integration (descent) & Adaptive 2nd-order Euler, exact apogee/ground root-solves & Same
3-DOF loop as ascent; termination at $z<0 \wedge v_z<0$ \\
\addlinespace
Motor mass & Manufacturer per-sample mass/CG when available; impulse-proportional fallback
for RASP & Constant-Isp 128-point reconstruction always; measured RSE mass/CG discarded \\
\addlinespace
Recovery & Parachute/streamer/tumble models, event-driven deployment with delays & Absent \\
\addlinespace
Validation & Flight-data validated: $\sim$10\% apogee subsonic; roll off $\sim 2\times$ &
Component-level analytic/table oracles ($10^{-4}$--$10^{-12}$); no whole-flight field data \\
\bottomrule
\end{tabularx}
\end{table}

Read column-wise, the table repeats the conclusion of \cref{app:or-features} in numerical
form: the gap is concentrated in what QtRocket \emph{consumes}, not in what it has built.
The one method where QtRocket is ahead --- a genuine embedded-pair integrator with step
rejection --- is also the one whose advantage will matter most on the day the dormant
machinery in the other rows is wired in.
